%-----------------------------------------------------------
%   Capítulo 2 - Estado da Arte 
%-----------------------------------------------------------
\chapter{Estado da Arte}
\label{cap:estado-arte}

\section{Comunicação em Sistemas Embebidos e Autónomos}

Nos sistemas embebidos e autónomos, as interfaces de comunicação respondem
a diferentes escalas de interligação e requisitos de funcionamento. O I2C
é utilizado na comunicação entre circuitos e periféricos, enquanto
interfaces como RS485 permitem interligar equipamentos em aplicações de
automação industrial e de
edifícios~\cite{nxp_um10204_2021,stiller_rs485_2015}.

Neste contexto, importa distinguir a camada física, relativa à transmissão
de sinais no meio, da camada de ligação de dados, que organiza a
transferência de tramas e o acesso ao meio. O protocolo de aplicação
define as operações e o significado dos dados trocados. A existência de
interfaces elétricas compatíveis não garante compatibilidade nas camadas
superiores~\cite{dodson_rs232_2023,stiller_rs485_2015,ieee_8023_2018}.

Estas tecnologias não ocupam todas o mesmo âmbito na comunicação.
Enquanto o I2C inclui regras próprias de transferência no barramento,
RS232 e RS485 caracterizam sobretudo a interface física. Modbus RTU
constitui um exemplo de protocolo que pode operar sobre ligações série,
sem ser uma propriedade intrínseca dessas
interfaces~\cite{nxp_um10204_2021,modbus_serial_2006}.

\section{I2C}

O I2C utiliza um barramento série síncrono partilhado constituído por duas
linhas: SDA, destinada à transferência de dados, e SCL, utilizada para o
sinal de relógio. Nos modos \textit{Standard-mode}, \textit{Fast-mode} e
\textit{Fast-mode Plus}, a sinalização baseia-se tipicamente em saídas de dreno ou
coletor aberto e na polarização das linhas por
\textit{pull-up}~\cite[sec.~3.1]{nxp_um10204_2021}.

Nestes modos, um controlador inicia cada transferência e gera o relógio;
o dispositivo alvo é selecionado por endereço de 7 ou 10 bits. Os dados circulam em bytes,
com indicação de confirmação ou não confirmação (ACK/NACK). A comunicação
é bidirecional, existindo arbitragem quando vários controladores disputam
o barramento~\cite[secs.~3.1.5--3.1.11]{nxp_um10204_2021}.

O reduzido número de linhas e o endereçamento integrado favorecem a
ligação local de periféricos, como sensores, memórias e expansores de
entradas e saídas, habitualmente no mesmo equipamento. Este contexto de
interligação entre circuitos distingue-se das redes distribuídas com
CAN ou RS485. A capacidade de expansão e a velocidade ficam condicionadas
pela capacitância do barramento e pelos tempos de subida dos
sinais~\cite[secs.~2 e~7]{nxp_um10204_2021}.

\section{RS232}

A interface RS232 estabelece ligações série ponto a ponto, com sinalização
não diferencial referida a uma massa comum. Os estados lógicos são
representados por tensões de polaridade positiva ou negativa; as linhas
distintas de transmissão e receção permitem comunicação simultânea nos
dois sentidos (\textit{full-duplex})~\cite[secs.~2.1 e~2.2]{dodson_rs232_2023}.

A norma define aspetos elétricos, funcionais e mecânicos. As designações
DTE (\textit{Data Terminal Equipment}) e DCE (\textit{Data Communication
Equipment}) distinguem funções dos equipamentos; deste âmbito não resulta
um mecanismo de endereçamento multiponto nem um protocolo de
aplicação~\cite[sec.~2]{dodson_rs232_2023}.
É frequente a utilização de comunicação assíncrona, mas a norma também
contempla sinais de temporização~\cite{ti_rs232_2002,dodson_rs232_2023}.

A RS232 permite uma ligação ponto a ponto simples e tem sido amplamente
utilizada na interligação de periféricos, modems e equipamentos de teste.
As principais limitações incluem a suscetibilidade ao ruído e as restrições de alcance e débito;
o comprimento admissível depende das características elétricas da
ligação, em particular da sua
capacitância~\cite{ti_rs232_2002,dodson_rs232_2023}.

\section{RS485}

RS485 define as características elétricas de emissores e recetores para
comunicação série diferencial multiponto. A informação é representada
pela diferença de tensão entre dois condutores, e vários nós podem
partilhar o barramento~\cite{stiller_rs485_2015,ti_rs485_an1057_2013}.

Na configuração \textit{half-duplex}, um par de sinais permite transmissão
alternada nos dois sentidos; a configuração \textit{full-duplex} utiliza
dois pares. A documentação recomenda uma disposição linear, com
derivações curtas, para reduzir problemas de
propagação~\cite[secs.~2 e~3]{ti_rs485_an1057_2013}.

A norma não define o protocolo de aplicação. Por conseguinte, o
endereçamento e a coordenação do acesso ao meio têm de ser tratados
acima da camada elétrica; a transmissão simultânea de vários emissores
no mesmo par pode causar contenção~\cite[secs.~1 e~10]{ti_rs485_an1057_2013}.

A robustez a ruído proporcionada pela sinalização diferencial e o suporte
multiponto constituem vantagens em aplicações de automação industrial,
automação de edifícios e controlo de motores~\cite{stiller_rs485_2015}.
O alcance e o débito são interdependentes, pelo que não se devem combinar
valores máximos obtidos em condições diferentes. A qualidade da comunicação
depende também da carga elétrica dos nós, das derivações e da adequação
das terminações para limitar reflexões~\cite[secs.~2 e~5--7]{ti_rs485_an1057_2013}.

\section{CAN}

CAN (\textit{Controller Area Network}) abrange mecanismos de ligação de
dados e opções de camada física. A família ISO~11898 distingue estes
âmbitos; CANopen e J1939 acrescentam protocolos e perfis superiores,
não sendo sinónimos de CAN~\cite{cia_can_standardization}.
Na camada física \textit{high-speed}, a sinalização é diferencial entre
CAN\_H e CAN\_L, recomendando-se um barramento linear terminado nas
extremidades. O alcance depende do débito e dos atrasos de propagação
da ligação~\cite{cia_can_hs}.

Em Classical CAN e CAN FD, as tramas são difundidas no barramento e identificadas por um campo de
11 ou 29 bits, associado ao conteúdo e à prioridade de acesso, em vez
de um endereço obrigatório de destinatário. Um protocolo superior pode
atribuir significado adicional a esse campo. Quando vários nós iniciam
a transmissão, a arbitragem compara os bits transmitidos: um nó que
envia um estado recessivo e observa um dominante cede o acesso, sem
destruir a trama vencedora~\cite{cia_can_cc}.

A deteção de erros inclui verificação por CRC, confirmação de receção
e monitorização dos bits. A deteção de um erro invalida a trama no nó
que o deteta e pode desencadear a sua retransmissão automática pelo
transmissor, de acordo com os mecanismos de tratamento de erros e o
estado do controlador; o confinamento de falhas pode colocar um nó em estado
\textit{bus-off}, impedindo que continue a transmitir. Estes mecanismos
e a arbitragem por prioridade favorecem o controlo embebido,
nomeadamente em veículos~\cite{cia_can_cc}.

Classical CAN admite até 8 bytes de dados por trama. CAN FD aumenta
esse limite para 64 bytes e permite, opcionalmente, um débito superior
na fase de dados, mantendo a arbitragem à velocidade nominal.
A dimensão das tramas e os condicionamentos físicos continuam a limitar
a transferência; a variante e a configuração da rede devem, por isso,
ser explicitadas na análise~\cite{cia_can_fd,cia_can_hs}.

\section{Ethernet}

IEEE~802.3 especifica a subcamada de controlo de acesso ao meio (MAC) e
as variantes de camada física (PHY) de Ethernet. O MAC organiza a
transferência de tramas; o PHY codifica e transmite os sinais no meio
selecionado. A norma contempla diferentes débitos, meios e comprimentos
de ligação, pelo que não existe um único alcance aplicável a todas
as variantes~\cite{ieee_8023_2018}.

Nas redes comutadas, os equipamentos estabelecem ligações ponto a ponto
a comutadores (\textit{switches}), formando habitualmente uma estrela
ou uma estrela hierárquica. As tramas contêm endereços MAC de origem e
destino e um campo FCS para deteção de erros. As ligações
\textit{full-duplex} permitem transmissão simultânea nos dois
sentidos~\cite{simmons_ethernet_2008}.

Ethernet pode transportar IP, sobre o qual operam TCP e UDP;
Modbus TCP é um exemplo de protocolo de aplicação neste contexto.
Estes protocolos pertencem a camadas superiores a
Ethernet~\cite{simmons_ethernet_2008,modbus_specifications}.
A diversidade de meios e débitos favorece a integração em redes de
monitorização e controlo. A adequação a sistemas embebidos e industriais
exige, contudo, delimitar a variante física, a organização da rede e
os serviços superiores necessários~\cite{ieee_8023_2018,simmons_ethernet_2008}.

A validação de requisitos com os docentes estabelece Ethernet como
ligação principal ao computador ou sistema central, com capacidade
nominal mínima de 100~Mbit/s; Gigabit Ethernet é desejável. A família
100BASE-T exemplifica Ethernet a 100~Mbit/s sobre pares entrançados;
o requisito fixa a capacidade nominal, mantendo a variante PHY e o meio
por selecionar. O débito nominal caracteriza a ligação,
enquanto o débito útil depende também dos campos de controlo, dos
intervalos entre tramas e dos protocolos
superiores~\cite{simmons_ethernet_2008}. Por conseguinte, esta capacidade
não garante uma latência ou um débito útil específicos do gateway, que
dependem ainda da carga e do processamento.

\section{Interface Física, Isolamento e Proteção}

A presença de uma função de comunicação digital não assegura, por si só,
os níveis elétricos exigidos pelo meio. RS232 requer emissores e recetores
adequados às suas tensões; RS485 requer circuitos de interface diferencial;
em CAN \textit{high-speed}, o transceiver realiza a ligação elétrica entre
o controlador e CAN\_H/CAN\_L. Estas funções podem estar integradas ou
exigir circuitos externos, conforme a plataforma, sem determinar nesta
fase componentes específicos~\cite{dodson_rs232_2023,ti_rs485_an1057_2013,cia_can_hs}.

Em ligações expostas ao ambiente industrial, a robustez depende também
das diferenças de potencial entre massas e da exposição a ruído,
descargas eletrostáticas (ESD) e sobretensões transitórias. A documentação
de RS485 apresenta isolamento galvânico para evitar caminhos de corrente
entre massas e proteção contra transientes dimensionada para as
perturbações esperadas; isolamento e proteção devem ser avaliados em
conjunto~\cite[secs.~1 e~6]{viviano_rs485_isolation_2018}.
No projeto, esta análise deve abranger a proteção da eletrónica de
controlo e a integridade da comunicação, de acordo com cada interface
e ambiente de utilização. Não decorre daqui a obrigatoriedade de isolar
todas as interfaces, nem a definição antecipada de níveis de imunidade.

\section{Sistemas ou Plataformas Multiprotocolo}

A literatura sobre \textit{gateways} multiprotocolo distingue a agregação
de tecnologias das funções de interpretação e troca de informação.
Abdelwahed et al.\ analisam componentes físicos e de software e
identificam desafios de compatibilidade, normalização e avaliação.
A revisão centra-se sobretudo em tecnologias sem fios; é aqui usada
para enquadrar a interoperabilidade, sem transpor resultados de
desempenho para I2C, RS485 ou CAN~\cite{abdelwahed_gateway_2025}.
% A consulta deste trabalho ficou limitada ao resumo, à introdução e aos
% excertos disponibilizados pela editora.

A análise seguinte seleciona quatro exemplos com contributos distintos:
duas plataformas programáveis e dois conversores para protocolos
definidos. Esta seleção permite comparar modelos de integração, sem
constituir uma escolha de plataforma para o projeto.

\subsection{Plataformas Embebidas Programáveis}

O GoWired Multiprotocol Gateway, da Domatic, exemplifica uma plataforma
baseada em microcontrolador, com software modificável e compatível com
o ambiente Arduino. A documentação anuncia duas interfaces RS485, CAN
e Ethernet, incluindo utilização simultânea de protocolos nas duas
interfaces série. Contudo, a página mantém a indicação de lançamento
futuro; estas características são tratadas como anunciadas, sem
confirmação de disponibilidade comercial ou de desempenho
experimental~\cite{domatic_gowired}.

O CompuLab IOT-GATE-iMX7 representa uma plataforma com suporte Linux
e interfaces configuráveis, incluindo uma ou duas portas Gigabit Ethernet
e PoE opcional~\cite{compulab_imx7}.
O seu manual evidencia uma limitação concreta: a opção CAN substitui
a interface RS485. I2C e SPI são disponibilizados no conector de
expansão P7; algumas funções alternativas partilham pinos, exigindo
verificação da configuração utilizada~\cite{compulab_imx7_manual_2018}.
A possibilidade de programar estas plataformas não demonstra, por si
só, a existência de conversão automática entre todas as interfaces.

\subsection{Gateways e Conversores Dedicados}

O Moxa MGate 5105-MB-EIP permite integrar Modbus RTU/ASCII/TCP e
EtherNet/IP. A documentação descreve a recolha de dados Modbus,
o armazenamento dos valores e a sua disponibilização em comunicações
EtherNet/IP, com mapeamento de dados em memória. A conversão tem,
assim, um âmbito definido pelos protocolos e funções suportados.
A porta série de dados é configurável como RS232, RS422 ou RS485,
não correspondendo a três portas independentes~\cite{moxa_mgate5105}.

O HMS Anybus X-gateway AB7838-F liga uma rede CANopen, onde assume
o papel de \textit{slave}, a EtherNet/IP, onde atua como
\textit{adapter}. O fabricante documenta a troca de dados cíclicos de
entradas e saídas e a configuração através de ferramenta própria.
Este exemplo distingue a integração entre protocolos superiores da
mera presença de CAN e Ethernet: as funções documentadas não
equivalem ao encaminhamento arbitrário de tramas~\cite{hms_ab7838}.

A Tabela~\ref{tab:plataformas-multiprotocolo} sintetiza apenas as
características relevantes para esta comparação. \enquote{NS} significa
não especificado nas fontes consultadas, e não ausência comprovada
da funcionalidade. \enquote{Integrada} designa uma interface disponibilizada
pelo equipamento; não implica que pertença diretamente ao MCU ou MPU.

\begin{table}[htbp]
    \centering
    \small
    \setlength{\tabcolsep}{4pt}
    \renewcommand{\arraystretch}{1.12}
    \caption{Plataformas programáveis e conversores dedicados selecionados.}
    \label{tab:plataformas-multiprotocolo}
    \begin{tabularx}{\textwidth}{@{}>{\raggedright\arraybackslash}p{0.18\textwidth}*{3}{>{\raggedright\arraybackslash}X}@{}}
        \toprule
        Sistema e fonte & Computação e software & Interfaces documentadas & Configuração e limites \\
        \midrule
        GoWired~\cite{domatic_gowired}
        & MCU ATSAMD21G18A; ambiente Arduino; SO/RTOS: NS.
        & Integradas anunciadas: duas RS485, CAN e Ethernet.
        RS232, CAN FD, I2C e SPI externos: NS.
        & Plataforma programável anunciada; simultaneidade série declarada,
        sem condições de desempenho publicadas. \\
        \addlinespace
        IOT-GATE-iMX7~\cite{compulab_imx7,compulab_imx7_manual_2018}
        & i.MX~7Dual; suporte Linux.
        & RS232 e RS485 integradas; CAN opcional em substituição de RS485;
        uma ou duas Gigabit Ethernet; I2C e SPI em P7. CAN FD: NS.
        & Opção CAN exclui RS485; funções alternativas em P7 partilham
        pinos. A integração depende do software. \\
        \addlinespace
        MGate 5105-MB-EIP~\cite{moxa_mgate5105}
        & MCU/MPU e SO/RTOS: NS.
        & Porta de dados RS232/422/485 e duas Ethernet integradas;
        RS232 de consola. CAN/CAN FD, I2C e SPI: NS.
        & Modbus--EtherNet/IP; mapeamento em memória.
        Modos elétricos alternativos na mesma porta de dados. \\
        \addlinespace
        Anybus AB7838-F~\cite{hms_ab7838}
        & MCU/MPU e SO/RTOS: NS.
        & CAN para CANopen e duas portas Ethernet integradas.
        RS232, RS485, CAN FD, I2C e SPI: NS.
        & CANopen \textit{slave}--EtherNet/IP \textit{adapter};
        troca de dados cíclicos de entradas e saídas. \\
        \bottomrule
    \end{tabularx}
\end{table}

A Tabela~\ref{tab:custos-plataformas} reúne a informação económica
disponível. \enquote{ND} significa informação de custo não determinada
nas fontes consultadas, sem implicar custo nulo.

\begin{table}[htbp]
    \centering
    \small
    \setlength{\tabcolsep}{4pt}
    \renewcommand{\arraystretch}{1.12}
    \caption{Informação económica disponível para as plataformas selecionadas.}
    \label{tab:custos-plataformas}
    \begin{tabularx}{\textwidth}{@{}>{\raggedright\arraybackslash}p{0.18\textwidth}*{2}{>{\raggedright\arraybackslash}X}@{}}
        \toprule
        Sistema & Custo de aquisição & Utilização / observações \\
        \midrule
        GoWired~\cite{domatic_gowired}
        & ND.
        & Utilização: ND; lançamento futuro anunciado. \\
        \addlinespace
        IOT-GATE-iMX7~\cite{compulab_imx7}
        & 129~USD por unidade; configuração mínima, encomendas de 1000 unidades.
        & Utilização: ND; não constitui cotação do protótipo. \\
        \addlinespace
        MGate 5105-MB-EIP~\cite{moxa_mgate5105}
        & ND.
        & Utilização: ND. \\
        \addlinespace
        Anybus AB7838-F~\cite{hms_ab7838}
        & ND.
        & Utilização: ND. \\
        \bottomrule
    \end{tabularx}
\end{table}

A página da CompuLab, consultada em 1 de outubro de 2026, publica o valor
de 129~USD por unidade para encomendas de 1000 unidades da configuração
mínima: 512~MB de RAM, 4~GB de eMMC e uma porta Ethernet. O fabricante
indica acréscimos por opções e fatores para outras quantidades, bem como
a possibilidade de alteração dos preços. Este valor não constitui uma
cotação para a configuração do futuro
protótipo~\cite{compulab_imx7}.

Nas páginas dos restantes produtos não foi determinado um preço de
aquisição. Também não foi determinado um custo total de utilização
comparável para qualquer das quatro soluções. A disponibilidade de
software modificável no GoWired ou de uma ferramenta de configuração
incluída no Anybus não permite concluir que a utilização seja
gratuita~\cite{domatic_gowired,hms_ab7838}.
A informação disponível limita a comparação económica: seria necessário
definir a configuração, os acessórios, o esforço de integração e
manutenção e as condições de operação para comparar custos totais.

\subsection{Normalização de Dados na Literatura Académica}

Castellanos et al.\ apresentam um \textit{gateway} que recebe informação
de nós sem fios e transforma os dados dos sensores num formato uniforme
para envio a um servidor. O estudo inclui uma prova de conceito com
dois nós, sustentando a discussão de normalização nesse cenário,
sem demonstrar conversão entre as cinco tecnologias deste
projeto~\cite{castellanos_gateway_2021}.

O interesse arquitetural deste caso reside na separação entre aquisição,
transformação dos dados e comunicação de saída. Por analogia, admite-se
como direção a avaliar uma organização com adaptadores por interface
ou protocolo e uma representação intermédia. Esta interpretação
não corresponde a uma arquitetura já validada para I2C, RS232, RS485,
CAN e Ethernet~\cite{castellanos_gateway_2021}.

\section{Concorrência e Adaptação de Taxas}

A utilização de várias interfaces em simultâneo, anunciada por exemplo
para as portas série do GoWired, coloca um problema de partilha dos
recursos de processamento e transmissão~\cite{domatic_gowired}.
Daí decorre a necessidade de gerir operações concorrentes e definir
como são atendidas quando disputam recursos. A arbitragem do próprio
barramento, como a prioridade dos identificadores CAN, não substitui
essa gestão no interior do gateway~\cite{cia_can_cc}. No projeto, a
política de prioridades terá de considerar os tempos de espera e evitar
que uma comunicação impeça indefinidamente o progresso das restantes;
esta necessidade não determina um sistema operativo ou escalonador.

As diferenças de taxa podem ocorrer tanto na concentração de dados dos
barramentos para Ethernet como no envio do sistema central para uma
interface mais lenta. Se os dados chegam antes de poderem ser
transmitidos, pode ser necessário armazenamento temporário
(\textit{buffering}). Por balanço entre a informação recebida e escoada,
um buffer finito absorve apenas desequilíbrios limitados: se a taxa de
entrada exceder de forma sustentada a capacidade de saída, não pode
garantir ausência de perdas. É necessário definir a carga admissível
e o comportamento perante saturação, em vez de atribuir essa garantia
apenas à existência de memória.

O controlo de fluxo ou \textit{backpressure} pode limitar a produção de
dados quando os protocolos e os participantes o permitem. Em TCP,
a janela anunciada pelo recetor condiciona os dados que o emissor pode
enviar e relaciona-se com o espaço de receção
disponível~\cite[sec.~3.8.6]{eddy_tcp_2022}. Este mecanismo exemplifica
uma possibilidade, sem assumir TCP no projeto nem a existência de
controlo de fluxo equivalente em todas as interfaces. As condições de
operação terão de compatibilizar taxas, armazenamento, prioridades e
tempos de espera, mantendo em aberto os limites quantitativos.

\section{Comparação das Tecnologias}

As Tabelas~\ref{tab:comparacao-organizacao}
e~\ref{tab:comparacao-condicoes} separam a organização da comunicação
dos condicionamentos de utilização. A comparação é qualitativa:
alcance e débito dependem da variante e da implementação, não sendo
apresentados máximos independentes como se fossem simultaneamente
atingíveis. Em CAN considera-se a camada física \textit{high-speed};
em Ethernet, a organização comutada, sem a generalizar a todas
as variantes.

\begin{table}[htbp]
    \centering
    \small
    \setlength{\tabcolsep}{4pt}
    \renewcommand{\arraystretch}{1.12}
    \caption{Organização da comunicação nas tecnologias estudadas.}
    \label{tab:comparacao-organizacao}
    \begin{tabularx}{\textwidth}{@{}>{\raggedright\arraybackslash}p{0.16\textwidth}*{3}{>{\raggedright\arraybackslash}X}@{}}
        \toprule
        Tecnologia & Topologia e multiponto & Sinalização & Identificação \\
        \midrule
        I2C~\cite{nxp_um10204_2021}
        & Barramento partilhado; vários dispositivos.
        & Série síncrona; SDA e SCL; não diferencial.
        & Endereços de 7 ou 10 bits. \\
        \addlinespace
        RS232~\cite{dodson_rs232_2023}
        & Ponto a ponto.
        & Não diferencial; transmissão e receção separadas.
        & Sem endereçamento multiponto próprio. \\
        \addlinespace
        RS485~\cite{ti_rs485_an1057_2013}
        & Barramento linear recomendado; multiponto.
        & Diferencial; um ou dois pares de sinais.
        & Definida pelo protocolo superior. \\
        \addlinespace
        CAN~\cite{cia_can_cc,cia_can_hs}
        & Barramento linear; vários nós.
        & Diferencial; estados dominante e recessivo.
        & Identificadores de trama; prioridade na arbitragem. \\
        \addlinespace
        Ethernet~\cite{simmons_ethernet_2008,ieee_8023_2018}
        & Ligações ponto a ponto a comutadores; rede com vários nós.
        & Dependente do PHY e do meio, elétrico ou ótico.
        & Endereços MAC de origem e destino. \\
        \bottomrule
    \end{tabularx}
\end{table}

\begin{table}[htbp]
    \centering
    \small
    \setlength{\tabcolsep}{4pt}
    \renewcommand{\arraystretch}{1.12}
    \caption{Condicionamentos e aplicações das tecnologias estudadas.}
    \label{tab:comparacao-condicoes}
    \begin{tabularx}{\textwidth}{@{}>{\raggedright\arraybackslash}p{0.16\textwidth}*{3}{>{\raggedright\arraybackslash}X}@{}}
        \toprule
        Tecnologia & Alcance e débito & Robustez e condicionantes & Aplicações típicas \\
        \midrule
        I2C~\cite{nxp_um10204_2021}
        & Condicionados pelo modo, capacitância e tempos de subida.
        & Dependente da carga e dos níveis elétricos.
        & Sensores, memórias e periféricos locais. \\
        \addlinespace
        RS232~\cite{ti_rs232_2002,dodson_rs232_2023}
        & Dependentes da capacitância e do cabo.
        & Suscetível a ruído; referência de massa comum.
        & Periféricos, modems e equipamentos de teste. \\
        \addlinespace
        RS485~\cite{ti_rs485_an1057_2013,stiller_rs485_2015}
        & Alcance e débito interdependentes.
        & Sinalização diferencial; atenção a terminações, derivações
        e carga.
        & Automação industrial e de edifícios. \\
        \addlinespace
        CAN~\cite{cia_can_hs,cia_can_cc,cia_can_fd}
        & Dependem da propagação e da variante; FD pode alterar
        o débito na fase de dados.
        & Deteção de erros e confinamento de falhas.
        & Controlo automóvel e sistemas distribuídos. \\
        \addlinespace
        Ethernet~\cite{ieee_8023_2018,simmons_ethernet_2008}
        & Definidos por variante, meio e condições da ligação.
        & FCS para deteção de erros; condições físicas dependentes do meio.
        & Redes de monitorização e controlo. \\
        \bottomrule
    \end{tabularx}
\end{table}

\section{Síntese Crítica do Estado da Arte}
\label{sec:sintese-estado-arte}

Os exemplos analisados evidenciam duas questões distintas:
disponibilizar interfaces e assegurar interoperabilidade. As plataformas
programáveis fornecem recursos para implementar a integração, sujeitos
às suas opções e incompatibilidades físicas; os conversores dedicados
documentam pares de protocolos, papéis e mapeamentos
concretos~\cite{compulab_imx7_manual_2018,moxa_mgate5105,hms_ab7838}.
Por conseguinte, o conceito de \textit{gateway} embebido multiprotocolo
em estudo deve ser delimitado pelas trocas de informação pretendidas,
sem pressupor conversão arbitrária entre quaisquer interfaces. No âmbito
confirmado na validação de requisitos com os docentes, Ethernet é a ligação principal ao
sistema central, que envia comandos e recebe dados de CAN, RS485, RS232
e I2C. O encaminhamento bidirecional entre Ethernet e cada uma destas
interfaces constitui o núcleo obrigatório; outros pares são uma
possibilidade de generalização a confirmar.

As tecnologias diferem na delimitação de tramas (\textit{framing}),
identificação, unidades de transferência, débito, temporização e
tratamento de erros. Por exemplo, RS485 define características elétricas,
não um tamanho máximo de dados de aplicação. A transmissão de grandes
volumes não é excluída por essa norma: a sua organização e viabilidade
dependem do protocolo superior e das condições da
ligação~\cite{ti_rs485_an1057_2013,modbus_serial_2006}.
Também não se pode pressupor correspondência de uma trama recebida
para uma trama transmitida. TCP ilustra esta diferença: transporta
um fluxo de bytes e não preserva os limites das operações de escrita
da aplicação~\cite[sec.~3.7]{eddy_tcp_2022}. A sua utilização neste exemplo
não constitui uma escolha de protocolo para o projeto.

Da comparação decorre que o encaminhamento pode exigir delimitação,
segmentação e adaptação, além da interpretação do protocolo superior.
O mapeamento
tem ainda de preservar o significado dos dados; copiar bytes não
estabelece, por si só, equivalência entre operações ou grandezas.
Esta conclusão articula os mecanismos de transporte com os exemplos
de mapeamento e normalização
estudados~\cite{eddy_tcp_2022,castellanos_gateway_2021}.

O acesso a partir do sistema central deve permitir selecionar a
interface e o destinatário ou identificador aplicável, fornecer dados
ou comandos e receber os resultados. A uniformidade e a transparência
deste acesso são objetivos desejáveis de abstração. Em CAN, o identificador
de trama não equivale necessariamente
a um endereço de nó; a sua interpretação depende do protocolo superior.
Esta abstração externa não obriga a uma representação intermédia
universal. Adaptadores e uma representação comum continuam a ser
hipóteses de projeto, a avaliar quanto à preservação da informação
específica e ao processamento adicional.

Para a Análise de Requisitos, estas conclusões sustentam o acesso
bidirecional centrado em Ethernet, a adaptação e preservação da informação
e a gestão de concorrência, prioridades e diferenças de taxa. A capacidade
geral de gateway exige informação suficiente sobre destino, configuração
e protocolo em cada operação; não implica conversão semântica automática
de protocolos desconhecidos. O mínimo nominal Ethernet de 100~Mbit/s
está definido, mas as condições de carga, os limites de armazenamento e
os critérios temporais e de fiabilidade permanecem por validar. A
observação de operações, erros e resultados deverá permitir verificar
esse comportamento sem antecipar a implementação.
